\pdfoutput=1
\documentclass[12pt,a4paper,reqno]{amsart}
\usepackage{amssymb}

\usepackage{amscd}
\usepackage{enumerate}
%\usepackage{psfig}
\usepackage{graphicx}
%\usepackage{showkeys}  
\usepackage{siunitx}
\usepackage{tikz-cd}
\usepackage{color}
\usetikzlibrary{arrows}
% uncomment this when editing cross-references
\numberwithin{equation}{section}

\usepackage{mathabx}


\usepackage{mathtools}%                  http://www.ctan.org/pkg/mathtools
\usepackage[tableposition=top]{caption}% http://www.ctan.org/pkg/caption
\usepackage{booktabs,dcolumn}%           http://www.ctan.org/pkg/dcolumn + http://www.ctan.org/pkg/booktabs

% Lighter notation.
\newcommand*\mc[1]{\multicolumn{1}{c}{#1}}
\newcommand*\tupref[2]{\href{http://math.mit.edu/~primegaps/tuples/admissible_#1_#2.txt}{\num{#2}}}

% Operators; uses ``mathtools''.
\DeclareMathOperator*\EH{EH}
\DeclareMathOperator*\GEH{GEH}
\DeclareMathOperator*\MPZ{MPZ}
\DeclareMathOperator*\DHL{DHL}
\DeclareMathOperator*\rad{rad}

%%\newcommand{\rad}[1]{{#1}^{\flat}}

%\DeclareMathOperator*\Kl{Kl} (commented because yield bad display for  \Kl_q %replaced with \newcommand... )
%\DeclareMathOperator*\FT{FT} (commented because yield bad display for \FT_q
%replaced with \newcommand...)

\DeclareMathOperator\swan{swan}
\DeclareMathOperator*\cond{cond}
\DeclareMathOperator*\Gal{Gal}

\newcommand{\FT}{\mathrm{FT}} 
\DeclareMathOperator{\hypk}{Kl}
\newcommand{\Kl}{\mathcal{K}\ell}
% Setup for ``caption''.
\DeclareCaptionLabelSeparator{separation}{:\quad}
\captionsetup{
  font=small,
  labelfont=sc,
  labelsep=separation,
  width=0.8\textwidth
}

\DeclareFontFamily{OT1}{rsfs}{}
\DeclareFontShape{OT1}{rsfs}{n}{it}{<-> rsfs10}{}
\DeclareMathAlphabet{\mathscr}{OT1}{rsfs}{n}{it}

\addtolength{\textwidth}{3 truecm}
\addtolength{\textheight}{1 truecm}
\setlength{\voffset}{-.6 truecm}
\setlength{\hoffset}{-1.3 truecm}
     
\theoremstyle{plain}

\newtheorem{theorem}{Theorem}[section]
%\newtheorem{theorem}[theorem]{Theorem}
\newtheorem{proposition}[theorem]{Proposition}
\newtheorem{lemma}[theorem]{Lemma}
\newtheorem{corollary}[theorem]{Corollary}
\newtheorem{conjecture}[theorem]{Conjecture}
\newtheorem{principle}[theorem]{Principle}
\newtheorem{claim}[theorem]{Claim}

\theoremstyle{definition}

%\newtheorem{roughdef}[subsection]{Rough Definition}
\newtheorem{definition}[theorem]{Definition}
\newtheorem{remark}[theorem]{Remark}
\newtheorem{remarks}[theorem]{Remarks}
\newtheorem{example}[theorem]{Example}
\newtheorem{examples}[theorem]{Examples}
%\newtheorem{problem}[subsection]{Problem}
%\newtheorem{question}[subsection]{Question}

\newcommand\F{\mathbb{F}}
\newcommand\E{\mathbb{E}}
\newcommand\R{\mathbb{R}}
\newcommand\Z{\mathbb{Z}}
\newcommand\N{\mathbb{N}}
\newcommand\C{\mathbb{C}}
\newcommand\Q{\mathbb{Q}}
\newcommand\eps{\varepsilon}
\newcommand\B{\operatorname{B}}
\renewcommand\Re{\operatorname{Re}}

\newcommand\Scal{\mathcal{S}}
\newcommand\Dcal[2]{\mathcal{D}^{({#1})}({#2})}
\newcommand\Dcone[1]{\mathcal{D}({#1})}
%%EK
\newcommand{\DI}[3]{\mathcal{D}_{{#1}}^{({#2})}({#3})}
\newcommand{\DIone}[2]{\mathcal{D}_{{#1}}({#2})}
\newcommand{\emt}{\mathrm{EMT}}
\newcommand{\uple}[1]{\text{\boldmath${#1}$}}
%%
%\newcommand\DHL{\operatorname{DHL}}
%\newcommand\EH{\operatorname{EH}}
%\newcommand\MPZ{\operatorname{MPZ}}
\newcommand{\opt}[1]{\textsf{{#1}}}
\newcommand{\algo}[1]{\textsf{{#1}}}

\newcommand\TypeI{\operatorname{Type}_{\operatorname{I}}}
\newcommand\TypeII{\operatorname{Type}_{\operatorname{II}}}
\newcommand\TypeIII{\operatorname{Type}_{\operatorname{III}}}
\renewcommand\P{\mathbb{P}}

\newcommand\rk{\mathrm{rk}}
\newcommand\Fr{\mathrm{Fr}}
\newcommand\tr{\mathrm{tr}}

\newcommand\ov{\overline}
\newcommand\wcheck{\widecheck}
%\newcommand\cond{\mathrm{cond}}
%\newcommand\Gal{\mathrm{Gal}}
%\newcommand\Kl{\mathrm{Kl}}
%\newcommand\swan{\mathrm{swan}}
%\newcommand\FT{\mathrm{FT}}
\newcommand\mcF{\mathcal{F}}
\newcommand\mcG{\mathcal{G}}
\newcommand\mcK{\mathcal{K}}
\newcommand\mcC{\mathcal{C}}
\newcommand\mcL{\mathcal{L}}
\newcommand\Gm{\mathbb{G}_m}
\newcommand\Af{\mathbb{A}^1}
\newcommand\Pl{\mathbb{P}^1}
\newcommand\Aa{\mathbb{A}}
\DeclareMathOperator{\Tr}{Tr}
\newcommand\Fp{\F_{p}}
\newcommand\Fpt{\F_{p}^{\times}}
\newcommand\Fq{\F_{q}}
\def\stacksum#1#2{{\stackrel{{\scriptstyle #1}}
{{\scriptstyle #2}}}}
\newcommand{\mods}[1]{\,(\mathrm{mod}\,{#1})}
\newcommand{\piar}{\pi_1}
\newcommand{\pige}{\pi_1^{g}}
\newcommand\sumsum{\mathop{\sum\sum}\limits}
\newcommand\multsum{\mathop{\sum\cdots\sum}\limits}
\newcommand\trpsum{\mathop{\sum\sum\sum}\limits}
%\newcommand\E{\mathbb{E}}
\newcommand\Rr{\mathbb{R}}
\newcommand\Zz{\mathbb{Z}}
\newcommand\Nn{\mathbb{N}}
\newcommand\Cc{\mathbb{C}}
\newcommand\Qq{\mathbb{Q}}
\newcommand{\what}{\widehat}
\newcommand{\ra}{\rightarrow}

\newcommand{\GL}{\mathrm{GL}}

%%%%%%%%%%%%%% EK %%%%%%%%%%%%%%
\newcommand{\ftd}{\mathcal{I}}
\renewcommand{\sim}{\asymp} 

%% This is more standard (for analytic number theorists) 
%% than \sim for A << B << A
\renewcommand{\lessapprox}{\llcurly}
\renewcommand{\gtrapprox}{\ggcurly}
%% I find this typographically better than \lessapprox
%% This symbol is from the mathabx package.
\renewcommand{\phi}{\varphi}
%% I think this form of phi is more standard for the Euler function
%% I think it is only used in improvements.tex for something else
\newcommand{\onef}{\mathbf{1}}
%% This is for characteristic functions
%%%%%%%%%%%%%%%%%%%%%%%%%%%%%%%%


\newcommand{\alert}[1]{{\bf \color{red} TODO: #1}}

%% Try normal paragraph indent
%%\parindent 0mm
%% Try to see what happens without enlarging the paragraph skip
%%\parskip   5mm 

%PhM use this for navigating the pdf when editing with texpad (comment if you prefer)
% (Ah, figured out why my LaTeX was not liking hyperref - there were old aux files that didn't use hyperref. -TT)
\usepackage{hyperref}

% for optimize.tex
\usepackage{algorithm, algorithmic}

\begin{document}

\title[Bounded intervals with many primes]{Variants of the Selberg sieve, and bounded intervals containing many primes}

\author{D.H.J. Polymath}
\address{http://michaelnielsen.org/polymath1/index.php}
%\email{}

\begin{abstract}  For any $m \geq 1$, let $H_m$ denote the quantity $\liminf_{n \to \infty} (p_{n+m}-p_n)$, where $p_n$ denotes the $n^{\operatorname{th}}$ prime.  A celebrated recent result of Zhang showed the finiteness of $H_1$, with the explicit bound $H_1 \leq \num{70000000}$, by combining a truncated version of the Selberg-type sieve of Goldston, Pintz, and Y{\i}ld{\i}r{\i}m with a new distributional estimate on primes.  This was then improved by us (the Polymath8 project) to $H_1 \leq \num{4680}$, and then by Maynard to $H_1 \leq \num{600}$, who also established for the first time a finiteness result for $H_m$ for $m \geq 2$, and specifically that $H_m \ll m^3 e^{4m}$.  Remarkably, these results of Maynard do not use any distributional results on the primes beyond the Bombieri-Vinogradov theorem, relying instead on a multidimensional version of the Selberg sieve.  If one also assumes the Elliott-Halberstam conjecture, Maynard obtained the bound $H_1 \leq 12$, improving upon the previous bound $H_1 \leq 16$ of Goldston, Pintz, and Y{\i}ld{\i}r{\i}m, as well as the bound $H_m \ll m^3 e^{2m}$.

In this paper, we extend the methods of Maynard by generalizing the Selberg sieve further, and by performing more extensive numerical calculations.  As a consequence, we can obtain the bound $H_1 \leq \num{246}$ unconditionally, and $H_1 \leq 6$ under the assumption of the generalized Elliott-Halberstam conjecture.  Indeed, under the latter conjecture we show the stronger statement that for any admissible triple $(h_1,h_2,h_3)$, there are infinitely many $n$ for which at least two of $n+h_1,n+h_2,n+h_3$ are prime.  We modify the ``parity problem'' argument of Selberg to show that the this result is the best possible that one can obtain one can obtain from purely sieve-theoretic considerations; in particular, no further improvement to the bound $H_1 \leq 6$ can be obtained by these methods.  For larger $m$, we use the distributional results obtained previously by our project to obtain the unconditional asymptotic bound $H_m \ll m e^{(4-\frac{24}{181})m}$, or $H_m \ll m e^{2m}$ under the assumption of the Elliott-Halberstam conjecture.  We also obtain explicit upper bounds for $H_m$ when $m=2,3,4,5$.
\end{abstract}

\subjclass[2010]{11P32}

\maketitle
%\today

\input introduction
\input subtheorems
\input sieving
\input more-sieving
\input asymptotics
\input h1-new
\input parity
\input remarks

\appendix
\input tuples
%\input exponential
%\input distribution
\input biblio


\end{document}